%Os materiais suplementares estão disponíveis no repositório público do projeto: 
Os códigos e dados utilizados neste estudo, disponíveis para uso não comercial, serão disponibilizados após a aceitação do artigo em: <<repositório GitHub do autor principal -- anonimizado devido ao processo double-blind review>>. 
\href{https://github.com/lucaszegrine/sus-prenatal-ai-icei-puc-minas-cc}%{github.com/lucaszegrine/sus-prenatal-ai-icei-puc-minas-cc}.

%O estudo não envolve participantes humanos, não utiliza dados clínicos reais identificáveis e, por isso, não foi submetido ao Comitê de Ética em Pesquisa da Universidade.
%PUC Minas. 

Este trabalho foi concebido com base em princípios de uso responsável de Inteligência Artificial, com especial atenção às diretrizes da Lei Geral de Proteção de Dados (Lei nº 13.709/2018 – LGPD). O sistema proposto adota processamento local (\textit{on-premise}), descarte imediato de áudio após transcrição e um \textit{gateway} de anonimização responsável por detectar e mascarar informações pessoalmente identificáveis (PII) antes da inferência pelo modelo, reduzindo a exposição de dados sensíveis.

O protótipo não utiliza dados clínicos reais nem envolve participantes humanos, baseando-se exclusivamente em cenários sintéticos para avaliação experimental. Dessa forma, o estudo não se enquadra como pesquisa envolvendo seres humanos segundo a Resolução CNS nº 510/2016, não sendo necessária submissão a Comitê de Ética em Pesquisa. Ainda assim, o desenho do sistema considera riscos potenciais associados ao contexto clínico, adotando mecanismos de mitigação como validação humana obrigatória (human-in-the-loop) e ausência de autonomia decisória do modelo.

Do ponto de vista ético, o sistema é explicitamente posicionado como ferramenta assistiva, não substituindo o julgamento clínico profissional. Para mitigar riscos de uso indevido ou respostas fora do escopo do pré-natal, são previstas salvaguardas adicionais, como restrições de domínio, possibilidade de recusa a consultas inadequadas e necessidade de validação manual antes da persistência de qualquer informação no prontuário.

Por fim, reconhece-se que, embora o protótipo demonstre viabilidade técnica, sua adoção em ambientes reais requer validação clínica, avaliação ética adicional e conformidade com regulamentações institucionais e profissionais aplicáveis.

Durante a preparação do texto, foi utilizado o modelo Gemini 3.1 Pro (Google, 04/2026) para revisão gramatical e refinamento de clareza.
A ferramenta não foi utilizada para geração de resultados experimentais, análise de resultados ou compartilhamento de dados sensíveis.
Todo o conteúdo obtido de forma automatizada foi revisado e editado conforme necessário, sob total responsabilidade dos autores.
